\chapter{When greedy is exact: the arc property}\label{sec:arc}

The counterexamples of \cref{sec:greedy} are thin polygons in which the
distance from a point is far from unimodal along the hull. This section shows
that this is the only obstruction: if the set of points far from $q$ is always
contiguous, the simplest greedy rule is exact, and for unimodal point sets it
gives a near-linear algorithm.

For $r>0$ and $q\in P$ let $\mathrm{Far}_r(q)=\{x\in P\setminus\{q\}:|qx|\ge r\}$.
We say that $P$ has the \emph{$r$-arc property} if for every $q$ the set
$\mathrm{Far}_r(q)$ consists of consecutive points in the order
$p_{q+1},\dots,p_{q-1}$; then $\mathrm{Far}_r(q)$ is exactly the ccw arc from
$a(q)$ to $b(q)$. We call $P$ \emph{unimodal} if for every $q$ the sequence
$|p_qp_{q+t}|$, $t=1,\dots,n-1$, is strictly increasing and then strictly
decreasing, where the maximum may be attained at two consecutive values of
$t$. A unimodal set has the $r$-arc property for every $r$. Points on a circle
are unimodal, since the chord length is $2R\sin(\theta/2)$ for the ccw angle
$\theta\in(0,2\pi)$; the tie at the top occurs, for instance, in every regular
polygon with an odd number of vertices, which is why it has to be admitted.

\begin{lemma}[Exchange]\label{lem:a1}
Assume the $r$-arc property. Let $S$ be an $r$-far set and let $s,s'$ be
consecutive points of $S$ in ccw order. Then $(S\setminus\{s'\})\cup\{a(s)\}$
is $r$-far.
\end{lemma}
\begin{proof}
By \cref{lem:arcbasic}, $a(s)$ lies on the ccw arc from $s$ to $s'$, and
$|s\,a(s)|\ge r$. Let $t\in S\setminus\{s,s'\}$. The set $\mathrm{Far}_r(t)$
is a set of consecutive points in the order that starts and ends next to $t$,
and it contains $s$ and $s'$. The ccw arc from $s$ to $s'$ does not contain
$t$, so all its points lie between $s$ and $s'$ in that order and belong to
$\mathrm{Far}_r(t)$. Hence $|t\,a(s)|\ge r$.
\end{proof}

For $q\in P$ and $k\ge2$ the \emph{greedy chain} of $q$ is $x_0=q$,
$x_{i+1}=a(x_i)$ for $i<k-1$. It is \emph{proper} if all $a(x_i)$ are defined
and the chain does not wrap around, i.e.\ the ccw offsets of $x_0,\dots,x_{k-1}$
from $q$ are strictly increasing and smaller than $n$.

\begin{theorem}\label{thm:arc}
Assume the $r$-arc property and let $k\ge2$.
\begin{enumerate}[label=(\roman*),itemsep=1pt,topsep=2pt]
\item If $S$ is an $r$-far $k$-set and $q\in S$, then the greedy chain of $q$
is proper and $r$-far.
\item A proper greedy chain is $r$-far if and only if $|x_0x_{k-1}|\ge r$.
\end{enumerate}
Hence $P$ has an $r$-far $k$-set if and only if some $q$ has a proper greedy
chain with $|x_0x_{k-1}|\ge r$.
\end{theorem}
\begin{proof}
(i) Let $S=\{s_1,\dots,s_k\}$ in ccw order with $s_1=q$. Applying
\cref{lem:a1} to the consecutive pair $(x_0,s_2)$ replaces $s_2$ by $x_1$;
then $(x_1,s_3)$ is a consecutive pair of the new $r$-far set, and we replace
$s_3$ by $x_2$, and so on. After $k-1$ steps the set is the chain, which is
$r$-far, and since $x_i$ lies on the ccw arc from $x_{i-1}$ to $s_{i+1}$, the
chain is proper.

(ii) Only ``if'' needs a proof. Fix $i$. Starting after $x_i$, the other
chain points appear in the order
$x_{i+1},\dots,x_{k-1},x_0,\dots,x_{i-1}$. The first of them is $a(x_i)$ and
lies in $\mathrm{Far}_r(x_i)$. The last one also does: for $i\ge1$ because
$x_i=a(x_{i-1})$, and for $i=0$ by assumption. Since $\mathrm{Far}_r(x_i)$ is
contiguous, it contains all of them.
\end{proof}

Part~(ii) is the difference to \cref{sec:greedy}: without the arc property a
greedy chain can close up and still contain a close pair, and by
\cref{thm:tightpair} there may be no $r$-far chain at all. Indeed the set
$P_\star$ violates the arc property at $r^2=10240$ ($\mathrm{Far}_r(p_0)$
consists of the three runs $\{p_1,p_2\}$, $\{p_4\}$, $\{p_6\}$).

\paragraph{Decision.}
Given the offsets of $a(q)$ for all $q$, the function $a^{k-1}$ together with
the offset travelled is obtained by repeated squaring in $O(n\log k)$ time and
$O(n)$ space, after which \cref{thm:arc} is tested in $O(1)$ per seed. (The
obvious $O(nk)$ walk works as well.) Computing $a(q)$ takes $O(n)$ time by
scanning. For unimodal $P$ we first find, for every row $q$, the position of
its maximum by binary search; then $a(q)$ is found by binary search on the
increasing part of the row.

\begin{theorem}\label{thm:arcdecision}
Under the $r$-arc property, the decision problem can be solved in
$O(n^2)$ time, and in $O(n\log n)$ time if $P$ is unimodal. A witness is
reported within the same bounds.
\end{theorem}

\paragraph{Optimisation.}
For unimodal $P$, sorting the $\binom n2$ distances and binary search gives
$O(n^2\log n)$ time; we call this variant \textsc{O1}. To avoid generating
all distances, observe that every row splits at its maximum into an increasing
and a decreasing list, so that the candidates are available as $2n$ implicit
sorted lists. Algorithm \textsc{O2} maintains an open interval $(lo,hi)$ with
$lo$ feasible and $hi$ infeasible and, for each list, the range of its entries
inside the interval. In each round it picks one of the remaining entries
uniformly at random, tests it, and shrinks all ranges by binary search.

\begin{theorem}\label{thm:arcopt}
For $n$ points in convex position that form a unimodal set, $k$-dispersion can
be solved exactly in expected $O(n\log^2n)$ time and $O(n)$ space.
\end{theorem}
\begin{proof}
A round costs $O(n\log n)$: $O(n)$ to select the pivot, $O(n\log n)$ for the
decision (\cref{thm:arcdecision}) and $O(n\log n)$ to update the ranges. All
remaining entries equal to the pivot are removed, and with probability at
least $1/2$ the pivot has rank between $1/4$ and $3/4$ among the remaining
entries, in which case at least a quarter of them are removed. So the expected
number of rounds is $O(\log n)$. When no entry is left, $lo=\rstar$, because
$\rstar$ is a pairwise distance and is feasible.
\end{proof}

\paragraph{Recognising the hypotheses.}
The arc property at a given $r$ is checked in $O(n^2)$ time by counting runs.
Unimodality can be verified faster.

\begin{lemma}\label{lem:unitest}
Whether $P$ is unimodal can be decided in $O(n\log n)$ time.
\end{lemma}
\begin{proof}
Fix $j$. The points strictly closer to $p_j$ than to $p_{j+1}$ are the points
of $P$ in an open half-plane bounded by the perpendicular bisector of
$p_jp_{j+1}$. Since $P$ is in strictly convex position, they form a run of
consecutive points; it contains $p_j$ and not $p_{j+1}$, so it is
$p_j,p_{j-1},\dots,p_{j-c_j+1}$ for some $1\le c_j\le n-1$. At most two points
of $P$ lie on the bisector; one boundary crossing of the bisector is inside the
edge $p_jp_{j+1}$, so the only candidate for a point equidistant from $p_j$ and
$p_{j+1}$ is $p_{j-c_j}$. Let $z_j=1$ if it is equidistant and $z_j=0$
otherwise. The predicate ``$p_{j-s}$ is strictly closer to $p_j$ than to
$p_{j+1}$'' is monotone in $s$, so $c_j$ is found by binary search with
$O(\log n)$ distance comparisons, and $z_j$ with one more.

Now fix a row $q$ and $1\le t\le n-2$, and put $j=q+t$. The step from
$p_{q+t}$ to $p_{q+t+1}$ increases the distance to $p_q$ iff $p_q=p_{j-t}$ is
strictly closer to $p_j$ than to $p_{j+1}$, that is, iff $t<c_j$; it leaves the
distance unchanged iff $t=c_j$ and $z_j=1$; otherwise it decreases. A row is
unimodal (with at most one tie, at the top) iff no step that is not an
increase is followed by a step that is not a decrease. Step $t$ is not an
increase iff $t\ge c_j$, and step $t+1$ (whose edge is $p_{j}p_{j+1}$ shifted
by one, i.e.\ index $j+1$) is not a decrease iff $t+1\le c_{j+1}-1+z_{j+1}$.
Hence some row fails iff for some $j$ there is an integer
$t\in[1,n-3]$ with $c_j\le t\le c_{j+1}+z_{j+1}-2$. As $c_j\ge1$, this is the
case iff $c_j\le\min\{c_{j+1}+z_{j+1}-2,\;n-3\}$ for some $j$, which is
checked in $O(n)$ time once all $c_j,z_j$ are known.
\end{proof}

For general convex input we use an \emph{adaptive} variant of \textsc{O1}: it
checks the arc property at every radius it tests and reports ``not
applicable'' if the check fails. It takes $O(n^2\log n)$ time, is exact
whenever it answers, and never answers incorrectly.

\paragraph{Weaker hypotheses.}
The hypothesis of \cref{thm:arc} cannot be weakened in the two obvious ways.

\begin{proposition}\label{prop:weaker}
\begin{enumerate}[label=(\alph*),itemsep=1pt,topsep=2pt]
\item There is a set of $10$ points and an $r$-far $6$-set $S$ such that
$\mathrm{Far}_r(s)$ is contiguous for every $s\in S$, but no greedy chain of six
points is $r$-far.
\item There is a set of $6$ points with the $\rstar$-arc property for $k=4$ on
which binary search over the candidate distances with the test of
\cref{thm:arc}(ii) returns a value larger than $\rstar$.
\end{enumerate}
\end{proposition}
\begin{proof}
The instances are listed in \code{arc/data/hypothesis_counterexamples.json} and
verified by exhaustive search in exact arithmetic
(\code{arc/check\_hypotheses.py}).
\end{proof}

Running times of the near-linear algorithm are reported in \cref{app:arcexp}.

